\section{Proofs of the estimates}\label{appBC}\label{appB}\label{app:subresultant}

This appendix proves the bounds that Section~\ref{sec5} states in $\poly$ form, recording the explicit
estimates. Notation is that of Sections~\ref{sec:prelim} and~\ref{sec5}; $\tau$ is the encoding length
of $(A,C,L_0,q)$; $\beta(x)$ is the largest bit length of the numerator and denominator of a rational
$x$, extended entrywise to matrices; logarithms are to base $2$; constants are not optimized.

Every field that occurs below is a number field in the representation of
Definition~\ref{sec5:def:model}: $K=\mathbb{Q}[\theta]/(f_K)$ with $f_K\in\mathbb{Z}[\theta]$
primitive of degree $d_K\le n$ and height $\beta_K=\poly(n,\tau)$, elements represented by coordinate
vectors in the power basis, and the \emph{size} of an element meaning the largest bit length of those
coordinates. Thus a statement such as ``height at most $\beta$'' always refers to that representation,
and every bound below is uniform in the field because $d_K$ and $\beta_K$ are themselves bounded by
$\poly(n,\tau)$ (Lemma~\ref{appBC:lem:bases}).

\subsection{Proof of Lemma~\ref{sec5:lem:height}: heights}\label{appBC:sub:norm}

Clear denominators in $[C;CA]$, multiplying by the least common multiple $D_0$ of the denominators:
this gives an integer matrix $\widehat W$ with $\beta(\widehat W)\le\tau+2pn\tau=:\beta_W$, because
$\beta(D_0)\le2pn\tau$. Its kernel is computed by Bareiss elimination, whose intermediate quantities
are minors, and Hadamard's inequality bounds a $j\times j$ minor by
$\log_2|\det|\le\frac{j}{2}\log_2 j+j\beta_W\le\frac{\rho}{2}\log_2\rho+\rho\beta_W$ for
$j\le\rho$. Clearing denominators in the resulting coordinates and dividing each column by the gcd of
its entries (the \emph{primitive normalization}) adds no bits and keeps $\det(N^{\mathsf T}N)$ an
integer, without which a common rational factor of $N$ would inflate both parts of every entry of $M$
\cite{hj13,vzgg03}.

\begin{lem}[explicit form of Lemma~\ref{sec5:lem:height}(a)]\label{appBC:lem:norm}\label{appBC:lem:comp}\label{appBC:cor:comp}
Let $\beta_W=\tau+2pn\tau$. Then
\begin{enumerate}[label=(\roman*),leftmargin=*]
\item $\beta(N)\le2\bigl(\frac{\rho}{2}\log_2\rho+\rho\beta_W\bigr)$, and $N$ is computable in bit
complexity $\poly(n,\tau)$;
\item $\beta(N^{\mathsf T}N)\le2\beta(N)+\log_2 n$ and
$\beta(N^{\mathsf T}AN)\le2\beta(N)+\tau+2\log_2 n$;
\item the entries of $M$ satisfy
$\beta(M)\le(k-1)\bigl(\tfrac12\log_2 k+\beta(N^{\mathsf T}N)\bigr)+\beta(N^{\mathsf T}AN)+\beta\bigl(\det(N^{\mathsf T}N)\bigr)$,
hence $\beta(M)\le c_1(n^{3}\tau+n^{2}\log_2 n)$ for a constant $c_1$.
\end{enumerate}
\end{lem}

\begin{proof}
(i) Each step has a polynomial arithmetic count in $n$ and acts on integers of $\poly(n,\tau)$ bits.
(ii) The entries of $N^{\mathsf T}N$ (respectively $N^{\mathsf T}AN$) are sums of $n$ (respectively
$n^{2}$) terms of size at most $2^{2\beta(N)}$ (respectively $2^{2\beta(N)+\tau}$). (iii) $N$ has full
column rank, so $N^{\mathsf T}N$ is integral and positive definite and $\det(N^{\mathsf T}N)\ge1$; the
numerator of $(N^{\mathsf T}N)^{-1}=\operatorname{adj}(N^{\mathsf T}N)/\det(N^{\mathsf T}N)$ is thus
controlled by the adjugate bound alone, whose entries are $(k-1)$-minors bounded by Hadamard
\cite{hj13,vzgg03}.
\end{proof}

\begin{lem}[explicit form of Lemma~\ref{sec5:lem:height}(b)--(d)]\label{appBC:lem:charpoly}\label{appBC:lem:root}\label{appBC:cor:root}\label{appBC:lem:bases}\label{appBC:sub:roots}
Let $K=\mathbb{Q}(\mu)$ be represented as in Definition~\ref{sec5:def:model}, with
$d_K=[K:\mathbb{Q}]\le n$ and $\beta_K=\beta(f_K)$.
\begin{enumerate}[label=(\roman*),leftmargin=*]
\item after clearing denominators and taking the primitive part,
$\beta(\chi_M)\le k\bigl(\beta(M)+\tfrac12\log_2 k\bigr)+\beta(M)+k=\poly(n,\tau)$;
\item every root $\mu$ of $\chi_M$ satisfies $d_K\le k\le n$ and
$h(\mu)\le\beta(\chi_M)+\tfrac12\log_2(k+1)+1$, so $\beta_K\le\beta(\chi_M)+\tfrac12\log_2(k+1)+1$;
the real roots are separable into pairwise disjoint rational intervals in bit complexity
$\poly(n,\tau)$ \cite{bpr06};
\item $E_\mu=\ker(M-\mu I)$ has over $K$ a primitive vector $a$ with
$\beta(a)\le2\bigl(\tfrac k2\log_2 k+k(\beta(M)+\beta_K)\bigr)$, bases of $E_\mu^0$,
$U_v^\perp$, $D_v=K_0\cap U_v^\perp$ and $\mathcal{Q}_v/\overline{V}_{0,v}$ are computable over $K$
with element sizes $\poly(n,\tau)$, and the bit complexity of all these computations is
$\poly(n,\tau)$.
\end{enumerate}
\end{lem}

\begin{proof}
(i) $\chi_M=\sum_i(-1)^{i}\sigma_i(M)t^{k-i}$ with $\sigma_i$ a sum of principal $i$-minors, each
bounded by $(\sqrt i\,\|M\|_\infty)^{i}$; summing $\binom ki\le2^{k}$ terms and clearing denominators
adds $\poly(n,\tau)$ bits, and Faddeev--LeVerrier or interpolation gives a polynomial arithmetic count.
(ii) Cauchy's bound gives $|\mu|\le1+\max_i|a_i/a_k|\le2^{\beta(\chi_M)+1}$, so each conjugate
satisfies $\log^{+}|\sigma(\mu)|\le\beta(\chi_M)+1$, and dividing by the degree gives the height bound
\cite{hj13}. (iii) The vector $a$ spans $\ker(M-\mu I)$ over $K$ and its coordinates are ratios of
$(k-1)$-minors over $K$ to a nonzero principal one, bounded by Hadamard together with
$\beta(M),\beta_K$; each of the four spaces is a kernel, image, sum or quotient of matrices over $K$
of size at most $n$, and all inner products lie in $K$; the complexity follows from $d_K\le n$, a
polynomial arithmetic count, and arithmetic in $K$ at cost $\poly(d_K,\beta_K)$ bits
\cite{lan02,vzgg03}.
\end{proof}

\begin{remark}[why the computation cannot be carried out over $\mathbb{Q}$]\label{appBC:rem:mustK}
When $\mu\notin\mathbb{Q}$ the space $E_\mu$ is a line over $K$ but not over $\mathbb{Q}$: its trace
over $\mathbb{Q}$ is a direct sum of $d$ conjugate lines. Computing the kernel over $\mathbb{Q}$ would
split one direction into $d$ conjugate directions and change both the fibre and the pencil, so
direction enumeration and fibre construction must take place in $K$. Since $\mu$ is never turned into a
number, nothing here grows with the distance from $\mu$ to $\mathbb{Q}$.
\end{remark}

\subsection{Proof of Lemma~\ref{sec5:lem:chart}: sizes of the charts}\label{appBC:sub:chart}

Fix a real eigenvalue $\mu$ and a direction $[v]\in\mathbb{P}(E_\mu^0)$; the fibre is
$\Gr(m,\mathcal{Q}_v/\overline{V}_{0,v})$ with $m=s-\delta$, and $\widetilde r:=r(v)-\delta$ is the
rank of the quotient bundle, where $r(v)=\rho$ for eigen directions and $\rho-1$ otherwise.

\begin{lem}[explicit form of Lemma~\ref{sec5:lem:chart}]\label{appBC:lem:chart}\label{appBC:rem:poly}
With $d_{\mathrm{chart}}=m(\widetilde r-m)$:
\begin{enumerate}[label=(\roman*),leftmargin=*]
\item $d_{\mathrm{chart}}\le\widetilde r^{2}/4\le p^{2}$;
\item an $\omega\times\omega$ maximal minor has degree at most $\omega$ in $X$ and in $\lambda$;
\item the number of minors is
$\binom q\omega\binom k\omega=k\binom qs\le n\sum_{j\le2p+1}\binom nj=n^{O(p)}$;
\item the chart cover of $\Gr(m,\widetilde r)$ has $\binom{\widetilde r}{m}\le2^{\widetilde r}\le2^{2p}$
elements, with at most $k\le n$ pieces, hence at most $n\cdot2^{2p}$ charts;
\item the dense size of a maximal minor, and of every subresultant below, is
\begin{equation}\label{appBC:eq:denseminor}
D=\binom{d_{\mathrm{chart}}+2\omega^{2}}{d_{\mathrm{chart}}},\qquad d_{\mathrm{chart}}\le p^{2},
\end{equation}
so $D\le\binom{p^{2}+2n^{2}}{p^{2}}=n^{O(p^{2})}$, a polynomial in $n$ once $p$ is fixed, the exponent
involving $p^{2}$ only.
\end{enumerate}
\end{lem}

\begin{proof}
$\widetilde r\le\rho\le2p$ and $x\mapsto x(\widetilde r-x)$ is maximal at $x=\widetilde r/2$, giving
(i); a chart row is affine in $X$ and the pencil is linear in $\lambda$, giving (ii);
$\binom k\omega=k$ because $\omega=k-1$, and $\binom qs\le\sum_{j\le2p}\binom nj$ because
$s\le\rho\le2p$, giving (iii); (iv) uses $\sum_j\binom{\widetilde r}j=2^{\widetilde r}$; (v) uses
$\binom{d+\omega}d\le\bigl(e(d+\omega)/d\bigr)^{d}$ with $d\le p^{2}$ constant and $\omega\le n$.
\end{proof}

A resultant of two $\lambda$-polynomials of degree $\omega$ is a Sylvester determinant of order
$2\omega$ with $X$-degree at most $\omega$ and height at most $\beta$, hence of bit length
$\poly(\omega,\beta)$; the criterion built on it is only \emph{sufficient}, failure being no certificate
that a chart is bad.

\subsection{Proof of Lemma~\ref{sec5:lem:subres}: the subresultant chain}\label{appBC:sub:subres}

\paragraph{Why naive Euclid fails}\label{appBC:sub:naive}
Division with remainder directly in $K(X)[\lambda]$ is the wrong procedure: from the second step on the
objects divided are rational functions, and each step multiplies the reciprocal of the leading
coefficient of the divisor into the next denominator, so denominator degrees can double at every step.
Take $K=\mathbb{Q}$ and one parameter $X$, with
\[
f=\lambda^{3}+X,\qquad g=\lambda^{2}+X\lambda+1 ,
\]
coprime over $\mathbb{Q}(X)[\lambda]$, and run the algorithm nevertheless. Dividing $f$ by $g$ gives
\[
q_1=\lambda-X,\qquad r_1=f-q_1g=(X^{2}-1)\lambda+2X ,
\]
with \emph{no} denominator; dividing $g$ by $r_1$ gives
\[
\begin{gathered}
q_2=\frac{\lambda}{X^{2}-1}+\frac{X^{3}-3X}{(X^{2}-1)^{2}},\\
r_2=g-q_2r_1=\frac{-X^{4}+4X^{2}+1}{(X^{2}-1)^{2}} ,
\end{gathered}
\]
a nonzero constant in $\lambda$, so $\gcd(f,g)=1$. The denominator jumps from none to degree $4$ in a
single step, and the numerator is exactly $\operatorname{Res}_\lambda(f,g)$; in general the quotient at
step $i$ carries $1/\operatorname{lc}_\lambda(r_{i-1})$, so denominators accumulate multiplicatively and
after $\omega$ steps can reach $2^{\Omega(\omega)}$. Direct Euclidean division over $K(X)[\lambda]$
therefore cannot bound coefficient size by $\poly(n,\tau)$.

\paragraph{The subresultant chain}
The correct tool is the \emph{subresultant}, a determinant computed inside the polynomial ring, so no
denominators arise. For $f=\sum_{i\le m}f_i\lambda^{m-i}$ and $g=\sum_{j\le n}g_j\lambda^{n-j}$ in
$K[X][\lambda]$, the Sylvester matrix $\operatorname{Syl}(f,g)$ is the $(m+n)\times(m+n)$ matrix of
shifted coefficient vectors with $\det=\operatorname{Res}(f,g)$, and for $0\le j<\min(m,n)$ the $j$-th
subresultant $S_j$ is the determinant obtained by deleting the last $j$ rows of each block and the last
$2j$ columns, so $S_0=\operatorname{Res}(f,g)$. Three properties matter: each $S_j$ lies in
$K[X][\lambda]$ (\emph{fraction-free}, the decisive advantage over Euclid), $\deg_\lambda S_j\le j$
(\emph{degree}), and with $j^{*}=\max\{j:S_j\neq0\}$ one has
$\deg_\lambda\gcd_{K(X)[\lambda]}(f,g)=j^{*}$ (\emph{gcd degree}), so $f,g$ are coprime if and only if
$S_0\neq0$. The sizes are as follows. The inputs satisfy
\begin{equation}\label{appBC:eq:sizesin}
\begin{gathered}
\deg_\lambda f,\ \deg_\lambda g\le\omega,\qquad
\deg_X f,\ \deg_X g\le\omega,\\
\text{coefficient size}\le\beta ,
\end{gathered}
\end{equation}
where $X$ carries $d\le p^{2}$ parameters and $\beta$ is measured in the representation of $K$. Hence
\begin{enumerate}[label=(\alph*),leftmargin=*]
\item $\operatorname{Syl}(f,g)$ has order $m+n\le2\omega$, and each of its entries is a coefficient of
$f$ or $g$, so of $X$-degree $\le\omega$ and size $\le\beta$;
\item the chain has $m+1\le\omega+1$ members $S_0,\dots,S_m$, and it is computed only up to the first
vanishing member, so its length is at most $\omega$;
\item each $S_j$ is a $(2\omega-j)\times(2\omega-j)$ determinant of such entries, hence
\begin{equation}\label{appBC:eq:sizesout}
\begin{gathered}
\deg_X S_j\le2\omega^{2},\\
\text{coefficient size}\le\beta_S,\\
\beta_S:=2\omega\bigl(\beta+\omega\log_2(2\omega)\bigr)+\poly(d_K,\beta_K),
\end{gathered}
\end{equation}
by Hadamard's inequality applied to the $2\omega$ rows;
\item the dense size of such an object is
\begin{equation}\label{appBC:eq:dense}
\begin{gathered}
D=\binom{d+2\omega^{2}}{d}=\binom{p^{2}+2\omega^{2}}{p^{2}},\\
D=n^{O(p^{2})}\ \text{for fixed } p,
\end{gathered}
\end{equation}
since $d\le p^{2}$ and $\omega\le n$; note that $D$ is polynomial in $n$ once $p$ is fixed, the
exponent depending on $p^{2}$ only.
\end{enumerate}
The cost of the fraction-free chain is $\poly(D,\beta_S)$, i.e.\ polynomial in the number of
coefficients and in their bit length, by (S4).

\begin{prop}[iteration and cost]\label{appBC:prop:iter}\label{appBC:prop:cost}\label{appBC:rem:pairwise}
Let $K$ be a number field in the representation of Definition~\ref{sec5:def:model}, and let
$m_1,\dots,m_\ell\in K[X][\lambda]$ satisfy \eqref{appBC:eq:sizesin} with $\ell=n^{O(p)}$. Then:
\begin{enumerate}[label=(\roman*),leftmargin=*]
\item $\gcd(m_1,\dots,m_\ell)=\gcd\bigl(\gcd(m_1,\dots,m_{\ell-1}),m_\ell\bigr)$, and iterating in the
order of the indices never increases the $\lambda$-degree or the $X$-degree;
\item the coefficient size of every iterated primitive gcd, and of every intermediate subresultant, is
at most $\beta_S$ of \eqref{appBC:eq:sizesout}, \emph{uniformly} over the iteration; it is bounded by
the input data, not decreased by the iteration;
\item the iterated algorithm computes the gcd within $\poly(\ell,D,\beta_S)$, and the test ``degree
positive'' agrees with ``all $\ell$ polynomials have a nontrivial common factor'';
\item a test based on \emph{pairwise} resultants would not do: for $f_1=\lambda(\lambda-1)$,
$f_2=\lambda(\lambda-2)$, $f_3=(\lambda-1)(\lambda-2)$ every pairwise resultant vanishes, yet
$\gcd(f_1,f_2,f_3)=1$.
\end{enumerate}
\end{prop}

\begin{proof}
(i) Associativity of the gcd; a gcd in $K[X][\lambda]$, taken primitive, divides its arguments in
$K[X][\lambda]$, so its $\lambda$-degree and its $X$-degree are bounded by those of the arguments and
never increase along the iteration. (ii) The primitive gcd of a step divides one of its arguments in
the UFD $K[X][\lambda]$ (Gauss's lemma), so by the factor-height bound (S5) its coefficient size is at
most the size bound attached to that argument, namely $\beta_S$ of \eqref{appBC:eq:sizesout}; the same
bound applies to every intermediate subresultant by item (a) of the description of the chain. The bound
is therefore uniform, but it is \emph{not} claimed to decrease from step to step: a gcd can be as
large, in coefficient size, as the input it divides. (iii) There are $\ell-1$ steps, each costing
$\poly(D,\beta_S)$ by (S4), and the total follows. (iv) Direct verification: $\lambda$ divides $f_1$,
$f_2$; $\lambda-1$ divides $f_1,f_3$; $\lambda-2$ divides $f_2,f_3$; no factor divides all three.
\end{proof}

\paragraph{Relation to the implementation}\label{appBC:sub:impl}
The complexity count of Section~\ref{sec5} is made for the subresultant chain, including its
multi-polynomial version; the implementation that produced the certificates of
Appendix~\ref{app:data} divides in $\mathbb{Q}(X)[\lambda]$ instead, which is feasible for the small
instances ($\omega\le3$, $d\le4$) with normalized denominators but costs what the example above
describes on large ones --- a limitation of the implementation, not of the theorem. Both compute the
\emph{same} object, the gcd in $K(X)[\lambda]$, whose degree-zero test is unaffected by removing the
content \cite{vzgg03}.

One further step passes from the gcd to a decision. Over a general family a positive gcd degree does
\emph{not} imply that every specialization has a common factor, since the factor may degenerate to a
constant ($X\lambda+1$ degenerates to $1$ at $X=0$); here the full-rank leading-coefficient lemma of
Section~\ref{sec4} supplies what is missing, because on a chart $\rank H\equiv\omega$, so at every
point some maximal minor has $\lambda$-degree exactly $\omega$ and the primitive common factor cannot
degenerate anywhere:
\[
\begin{gathered}
\bigl(\deg_\lambda\gcd>0\bigr)\ \wedge\ \bigl(\rank H\equiv\omega\bigr)\\
\Longrightarrow\quad \text{every point of the chart is bad}.
\end{gathered}
\]
\emph{Computing} the gcd (this appendix) and \emph{using} it to judge a whole chart
(Section~\ref{sec4}) are thus two separate matters.


